% Generated by stress/numerics/report/make_numbers_extra.py -- do not edit.
% Computed from the full-run results/*.jsonl (and summary.json for e04 convergence).

% source: e01 sweep numpy: batch size N of the peak throughput
\newcommand{\NumMapPeakNumpyN}{\ensuremath{10^{4}}}
% source: e01 sweep numpy: throughput at the largest N (points/s)
\newcommand{\NumMapPlateauNumpy}{\ensuremath{1.7\times10^{8}}}
% source: e01 sweep numpy: peak throughput / throughput at the largest N
\newcommand{\NumMapPeakOverPlateauNumpy}{\ensuremath{5.4\times}}
% source: e01 sweep: map_batch GPU / map_batch CPU throughput at the largest N
\newcommand{\NumMapBatchGpuOverCpuMaxN}{\ensuremath{1.11\times}}
% source: e01 sweep: cupy end-to-end / numpy throughput at the largest N
\newcommand{\NumRawGpuOverCpuMaxN}{\ensuremath{0.87\times}}
% source: e02 profile k28 cpu: refinement share of tottime
\newcommand{\NumProfRefinementKTwoEight}{46\%}
% source: e02 profile k28 cpu: linked-list share of tottime
\newcommand{\NumProfLinkedListKTwoEight}{38\%}
% source: e02 profile k28 cpu: other share of tottime
\newcommand{\NumProfOtherKTwoEight}{17\%}
% source: e02 profile k28 cpu: total profiled seconds
\newcommand{\NumProfTotalSKTwoEight}{83}
% source: e02 profile k28 cpu: ManifoldMachine._refine_layer tottime / total
\newcommand{\NumProfRefineLayerShareKTwoEight}{44\%}
% source: e02 profile k28 cpu: Point.__init__ calls
\newcommand{\NumProfPointInitCallsKTwoEight}{\ensuremath{8.7\times10^{6}}}
% source: e02 profile k28 cpu: Point.__init__ cumtime / total
\newcommand{\NumProfPointInitShareKTwoEight}{18\%}
% source: e02 profile k28 cpu: numpy.array calls
\newcommand{\NumProfNpArrayCallsKTwoEight}{\ensuremath{8.7\times10^{6}}}
% source: e02 profile k28 cpu: refinement + linked-list share of tottime
\newcommand{\NumProfRefPlusListKTwoEight}{83\%}
% source: e02 profile p3 cpu: refinement share of tottime
\newcommand{\NumProfRefinementPThree}{30\%}
% source: e02 profile p3 cpu: linked-list share of tottime
\newcommand{\NumProfLinkedListPThree}{45\%}
% source: e02 profile p3 cpu: other share of tottime
\newcommand{\NumProfOtherPThree}{25\%}
% source: e02 profile p3 cpu: refinement + linked-list share of tottime
\newcommand{\NumProfRefPlusListPThree}{75\%}
% source: e02 profile k10 cpu: refinement share of tottime
\newcommand{\NumProfRefinementKTen}{42\%}
% source: e02 profile k10 cpu: linked-list share of tottime
\newcommand{\NumProfLinkedListKTen}{33\%}
% source: e02 profile k10 cpu: other share of tottime
\newcommand{\NumProfOtherKTen}{24\%}
% source: e02 profile k10 cpu: refinement + linked-list share of tottime
\newcommand{\NumProfRefPlusListKTen}{76\%}
% source: e02 verify k28: CPU unstable points at the last step (12)
\newcommand{\NumGpuVerifyPointsCpuKTwoEight}{4\,348\,557}
% source: e02 verify k28: GPU unstable points at the last step (12)
\newcommand{\NumGpuVerifyPointsGpuKTwoEight}{4\,348\,603}
% source: e02 verify k28: max |coord| at the last step
\newcommand{\NumGpuVerifyMaxCoordKTwoEight}{\ensuremath{5.5\times10^{8}}}
% source: e02 verify k28: first step with a non-zero CPU-GPU deviation
\newcommand{\NumGpuFirstDivergentStepKTwoEight}{8}
% source: e02 verify k28: max |CPU-GPU| coordinate at step 11
\newcommand{\NumGpuVerifyDevStepElevenKTwoEight}{\ensuremath{3.6\times10^{-11}}}
% source: e02 verify k28: max |CPU-GPU| / max|coord| at step 11
\newcommand{\NumGpuVerifyRelDevStepElevenKTwoEight}{\ensuremath{1.5\times10^{-15}}}
% source: e02 timed k28: unstable points after the last growth step
\newcommand{\NumGpuTimedMaxPointsKTwoEight}{\ensuremath{4.3\times10^{6}}}
% source: e03 growth k=2.8 cutoff 1e-07: point growth factor of the last completed step
\newcommand{\NumLastGrowthFactorKTwoEight}{24.0}
% source: e03 growth k=2.8: unstable eigenvalue lambda_u
\newcommand{\NumLambdaUKTwoEight}{5.72}
% source: e03 growth k=2.8 cutoff 1e-07: max|coord| after the last completed step
\newcommand{\NumLastMaxCoordKTwoEight}{\ensuremath{1.3\times10^{6}}}
% source: e03 growth k=2.8 cutoff 1e-07: unstable points after the last completed step
\newcommand{\NumLastPointsKTwoEight}{407\,709}
% source: e03 growth k=2.8 cutoff 1e-07: completed growth steps
\newcommand{\NumGrowthStepsKTwoEight}{12}
% source: e03 growth k=2.8 cutoff 1e-05: point growth factor of the last completed step
\newcommand{\NumLastGrowthFactorKTwoEightCoarse}{23.4}
% source: e03 growth k=2.8: unstable eigenvalue lambda_u
\newcommand{\NumLambdaUKTwoEightCoarse}{5.72}
% source: e03 growth k=2.8 cutoff 1e-05: max|coord| after the last completed step
\newcommand{\NumLastMaxCoordKTwoEightCoarse}{\ensuremath{1.3\times10^{6}}}
% source: e03 growth k=2.8 cutoff 1e-05: unstable points after the last completed step
\newcommand{\NumLastPointsKTwoEightCoarse}{86\,168}
% source: e03 growth k=2.8 cutoff 1e-05: completed growth steps
\newcommand{\NumGrowthStepsKTwoEightCoarse}{12}
% source: e03 growth k=4.0 cutoff 1e-07: point growth factor of the last completed step
\newcommand{\NumLastGrowthFactorKFour}{12.9}
% source: e03 growth k=4.0: unstable eigenvalue lambda_u
\newcommand{\NumLambdaUKFour}{6.31}
% source: e03 growth k=4.0 cutoff 1e-07: max|coord| after the last completed step
\newcommand{\NumLastMaxCoordKFour}{\ensuremath{1.6\times10^{5}}}
% source: e03 growth k=4.0 cutoff 1e-07: unstable points after the last completed step
\newcommand{\NumLastPointsKFour}{89\,077}
% source: e03 growth k=4.0 cutoff 1e-07: completed growth steps
\newcommand{\NumGrowthStepsKFour}{11}
% source: e03 growth k=6.0 cutoff 1e-07: point growth factor of the last completed step
\newcommand{\NumLastGrowthFactorKSix}{61.8}
% source: e03 growth k=6.0: unstable eigenvalue lambda_u
\newcommand{\NumLambdaUKSix}{7.15}
% source: e03 growth k=6.0 cutoff 1e-07: max|coord| after the last completed step
\newcommand{\NumLastMaxCoordKSix}{\ensuremath{3.5\times10^{7}}}
% source: e03 growth k=6.0 cutoff 1e-07: unstable points after the last completed step
\newcommand{\NumLastPointsKSix}{2\,194\,304}
% source: e03 growth k=6.0 cutoff 1e-07: completed growth steps
\newcommand{\NumGrowthStepsKSix}{11}
% source: e03 growth k=10.0 cutoff 1e-07: point growth factor of the last completed step
\newcommand{\NumLastGrowthFactorKTen}{38.6}
% source: e03 growth k=10.0: unstable eigenvalue lambda_u
\newcommand{\NumLambdaUKTen}{8.52}
% source: e03 growth k=10.0 cutoff 1e-07: max|coord| after the last completed step
\newcommand{\NumLastMaxCoordKTen}{\ensuremath{6.6\times10^{6}}}
% source: e03 growth k=10.0 cutoff 1e-07: unstable points after the last completed step
\newcommand{\NumLastPointsKTen}{520\,502}
% source: e03 growth k=10.0 cutoff 1e-07: completed growth steps
\newcommand{\NumGrowthStepsKTen}{10}
% source: e03 growth: configs killed at the 24 GB RSS cap
\newcommand{\NumScalingOomConfigs}{3}
% source: e03 growth oom: fewest unstable points held after the last completed step
\newcommand{\NumScalingOomMinPointsHeld}{86\,168}
% source: e03 growth oom: most unstable points held after the last completed step
\newcommand{\NumScalingOomMaxPointsHeld}{407\,709}
% source: e03 growth oom: process wall minus checkpointed wall = time spent in the killed step (min, s)
\newcommand{\NumScalingOomKilledStepMinS}{190}
% source: e03 growth oom: time spent in the killed step (max, s)
\newcommand{\NumScalingOomKilledStepMaxS}{208}
% source: derived: 24 GiB / median RSS slope 1057.7 B/point (summary.json BytesPerPoint)
\newcommand{\NumPointCeilingTwentyFourGB}{\ensuremath{2.4\times10^{7}}}
% source: e04 ('k28', 10): log-log slope of the p50 invariance error vs cutoff, cutoffs 1e-9..1e-6
\newcommand{\NumInvSlopeKTwoEightTen}{0.70}
% source: e04 ('k28', 10): crossings at cutoff 1e-9
\newcommand{\NumCrossingsKTwoEightTenRef}{5}
% source: e04 ('k28', 10): crossings at cutoff 1e-3
\newcommand{\NumCrossingsKTwoEightTenCutoffThree}{5}
% source: e04 ('k28', 10): crossings at cutoff 1e-4
\newcommand{\NumCrossingsKTwoEightTenCutoffFour}{5}
% source: e04 ('k28', 10): crossings at cutoff 1e-5
\newcommand{\NumCrossingsKTwoEightTenCutoffFive}{5}
% source: e04 ('k28', 12): log-log slope of the p50 invariance error vs cutoff, cutoffs 1e-9..1e-6
\newcommand{\NumInvSlopeKTwoEightTwelve}{0.70}
% source: e04 ('k28', 12): crossings at cutoff 1e-9
\newcommand{\NumCrossingsKTwoEightTwelveRef}{17}
% source: e04 ('k28', 12): crossings at cutoff 1e-3
\newcommand{\NumCrossingsKTwoEightTwelveCutoffThree}{17}
% source: e04 ('k28', 12): crossings at cutoff 1e-4
\newcommand{\NumCrossingsKTwoEightTwelveCutoffFour}{17}
% source: e04 ('k28', 12): crossings at cutoff 1e-5
\newcommand{\NumCrossingsKTwoEightTwelveCutoffFive}{17}
% source: e04 ('p3', 13): log-log slope of the p50 invariance error vs cutoff, cutoffs 1e-9..1e-6
\newcommand{\NumInvSlopePThreeThirteen}{0.64}
% source: e04 ('p3', 13): crossings at cutoff 1e-9
\newcommand{\NumCrossingsPThreeThirteenRef}{27}
% source: e04 ('p3', 13): crossings at cutoff 1e-3
\newcommand{\NumCrossingsPThreeThirteenCutoffThree}{14}
% source: e04 ('p3', 13): crossings at cutoff 1e-4
\newcommand{\NumCrossingsPThreeThirteenCutoffFour}{37}
% source: e04 ('p3', 13): crossings at cutoff 1e-5
\newcommand{\NumCrossingsPThreeThirteenCutoffFive}{27}
% source: e04 ('p3', 17): log-log slope of the p50 invariance error vs cutoff, cutoffs 1e-9..1e-6
\newcommand{\NumInvSlopePThreeSeventeen}{0.66}
% source: e04 ('p3', 17): crossings at cutoff 1e-9
\newcommand{\NumCrossingsPThreeSeventeenRef}{203}
% source: e04 ('p3', 17): crossings at cutoff 1e-4
\newcommand{\NumCrossingsPThreeSeventeenCutoffFour}{249}
% source: e04 ('p3', 17): crossings at cutoff 1e-5
\newcommand{\NumCrossingsPThreeSeventeenCutoffFive}{214}
% source: e04 k28/10: log-log slope of the median crossing distance to the 1e-9 run vs cutoff, cutoffs <= 1e-5
\newcommand{\NumCrossDevSlopeKTwoEightTen}{0.73}
% source: e04 k28/12: log-log slope of the median crossing distance to the 1e-9 run vs cutoff, cutoffs <= 1e-5
\newcommand{\NumCrossDevSlopeKTwoEightTwelve}{0.70}
% source: e04 p3/13: log-log slope of the median crossing distance to the 1e-9 run vs cutoff, cutoffs <= 1e-5
\newcommand{\NumCrossDevSlopePThreeThirteen}{0.66}
% source: e04 p3/17: log-log slope of the median crossing distance to the 1e-9 run vs cutoff, cutoffs <= 1e-5
\newcommand{\NumCrossDevSlopePThreeSeventeen}{0.58}
% source: e05 k10: ||lu ls| - 1|
\newcommand{\NumDetDefectKTen}{\ensuremath{1.2\times10^{-11}}}
% source: e05 k10: max |stored - analytic| Jacobian entry (the k10 recipe passes no analytic Jacobian)
\newcommand{\NumJacDevKTen}{\ensuremath{1.4\times10^{-12}}}
% source: e05 k28 two blasts: |signed area| of the turnstile-sized lobes
\newcommand{\NumLobeAreaKTwoEight}{7.004}
% source: e05 k28 two blasts: lobes of |area| within 0.1 of 7
\newcommand{\NumLobeAreaKTwoEightCount}{2}
% source: e05 k28 two blasts: max |A_image - A| over those lobes
\newcommand{\NumLobeAreaAbsErrKTwoEight}{\ensuremath{5.3\times10^{-6}}}
% source: e05 p3: median |signed lobe area|
\newcommand{\NumLobeAreaPThreeMedian}{\ensuremath{1.1\times10^{-4}}}
% source: e05 p3: max |A_image - A| over the lobes
\newcommand{\NumLobeAbsErrPThreeMax}{\ensuremath{5.4\times10^{-7}}}
% source: e05 p3: lobe pairs compared
\newcommand{\NumLobePairsPThree}{22}
% source: e05: max ||lu ls| - 1| over every case except k10 (analytic Jacobian)
\newcommand{\NumDetDefectMaxAnalytic}{\ensuremath{4.4\times10^{-16}}}
% source: e06 b=1 saddle: 'Max iterations reached' exceptions
\newcommand{\NumSweepMaxIterFails}{10}
% source: e06: smallest k of those
\newcommand{\NumSweepMaxIterKMin}{0.30}
% source: e06: largest k of those
\newcommand{\NumSweepMaxIterKMax}{0.545}
% source: e06 b=1 saddle: configs ok through trellis_pips
\newcommand{\NumSweepSaddleOk}{50}
% source: e06 b=1 saddle ok: configs with 0 non-anchor crossings
\newcommand{\NumSweepZeroCrossings}{26}
% source: e06 b=1 saddle ok: largest k with 0 crossings
\newcommand{\NumSweepZeroCrossingsKMax}{5.93}
% source: e06 b=1 saddle ok: smallest k with >= 1 crossing
\newcommand{\NumSweepFirstCrossingK}{2.50}
% source: e06 b=1 saddle ok: most crossings in any config
\newcommand{\NumSweepMaxCrossings}{3}
% source: e06 b=1 other point: 'not a saddle' exceptions
\newcommand{\NumSweepOtherElliptic}{6}
% source: e06: largest k of those
\newcommand{\NumSweepOtherEllipticKMax}{2.63}
% source: e06 b=1 other point: oom at the 4 GB cap
\newcommand{\NumSweepOtherOom}{3}
% source: e06: smallest k of those
\newcommand{\NumSweepOtherOomKMin}{10.2}
% source: e06 b=-1: exceptions
\newcommand{\NumSweepBMinusOneFails}{8}
% source: e06 b in {0.3,0.5,0.9}: ok / configs
\newcommand{\NumSweepDissipativeOk}{15/15}
% source: e07 blasts p3: crossings before blasting
\newcommand{\NumBlastCrossingsStartPThree}{30}
% source: e07 blasts p3: crossings after the last blast
\newcommand{\NumBlastCrossingsEndPThree}{75}
% source: e07 blasts p3: time of the last blast (s)
\newcommand{\NumBlastLastSPThree}{0.0093}
% source: e07 blasts p3: bridge classes after the last blast
\newcommand{\NumBlastClassesPThree}{31}
% source: e07 blasts p3: unresolved classes
\newcommand{\NumBlastUnresolvedPThree}{1}
% source: e07 blasts p3: virtual classes
\newcommand{\NumBlastVirtualPThree}{1}
% source: e07 blasts nested: crossings before blasting
\newcommand{\NumBlastCrossingsStartNested}{34}
% source: e07 blasts nested: crossings after the last blast
\newcommand{\NumBlastCrossingsEndNested}{276}
% source: e07 blasts nested: time of the last blast (s)
\newcommand{\NumBlastLastSNested}{1.5}
% source: e07 blasts nested: bridge classes after the last blast
\newcommand{\NumBlastClassesNested}{13}
% source: e07 blasts nested: unresolved classes
\newcommand{\NumBlastUnresolvedNested}{0}
% source: e07 blasts nested: virtual classes
\newcommand{\NumBlastVirtualNested}{0}
% source: e07 blasts k28: crossings before blasting
\newcommand{\NumBlastCrossingsStartKTwoEight}{4}
% source: e07 blasts k28: crossings after the last blast
\newcommand{\NumBlastCrossingsEndKTwoEight}{712}
% source: e07 blasts k28: time of the last blast (s)
\newcommand{\NumBlastLastSKTwoEight}{5.6}
% source: e07 blasts k28: bridge classes after the last blast
\newcommand{\NumBlastClassesKTwoEight}{32}
% source: e07 blasts k28: unresolved classes
\newcommand{\NumBlastUnresolvedKTwoEight}{7}
% source: e07 blasts k28: virtual classes
\newcommand{\NumBlastVirtualKTwoEight}{2}
% source: e07 escape fixed: process wall minus checkpointed wall (time inside the killed step, s)
\newcommand{\NumEscapeKilledStepS}{175}
% source: e07 depth p3 18: unstable points after step 17, before the kill
\newcommand{\NumPThreePointsSeventeen}{584\,747}
% source: e07 depth p3 18: time of growth step 17 (s)
\newcommand{\NumPThreeStepSeventeenS}{7.2}
% source: e07 depth inversion 7: unstable points after step 6, before the kill
\newcommand{\NumInversionPointsSixSteps}{49\,417}
% source: e07 depth p3: deepest run with is_reliable = True
\newcommand{\NumDepthReliableMaxPThree}{14}
% source: e07 depth p3: shallowest run with is_reliable = False
\newcommand{\NumDepthUnreliableMinPThree}{15}
% source: e07 depth p3: unstable steps whose run hit the RSS cap
\newcommand{\NumDepthOomPThree}{18}
% source: e07 depth inversion: unstable steps whose topology stack raised an invariant AssertionError
\newcommand{\NumInversionAssertSteps}{4 and 5}
% source: e07 depth inversion: unstable steps whose run hit the RSS cap
\newcommand{\NumInversionOomSteps}{7 and 8}
% source: e07 tangency: configs (24 k values x 3 cutoffs)
\newcommand{\NumTangencyConfigs}{72}
% source: e07 tangency k+1=1e-4 cutoff 1e-07: non-anchor crossings
\newcommand{\NumTangencyCrossingsNearCutoffSeven}{25}
% source: e07 tangency cutoff 1e-07: most non-anchor crossings over k
\newcommand{\NumTangencyCrossingsPeakCutoffSeven}{149}
% source: e07 tangency k+1=1e-4 cutoff 1e-10: non-anchor crossings
\newcommand{\NumTangencyCrossingsNearCutoffTen}{153}
% source: e07 tangency cutoff 1e-10: most non-anchor crossings over k
\newcommand{\NumTangencyCrossingsPeakCutoffTen}{1\,607}
% source: e07 tangency k+1=1e-4 cutoff 1e-12: non-anchor crossings
\newcommand{\NumTangencyCrossingsNearCutoffTwelve}{499}
% source: e07 tangency cutoff 1e-12: most non-anchor crossings over k
\newcommand{\NumTangencyCrossingsPeakCutoffTwelve}{4\,401}
% source: e07 tangency: smallest k+1 from which all three cutoffs give the same crossing count
\newcommand{\NumTangencyAgreeKPlusOne}{0.16}
% source: e07 tangency k+1 = 1: non-anchor crossings (every cutoff)
\newcommand{\NumTangencyCrossingsFar}{7}
% source: e07 tangency k+1=1e-4: heuristic splitting reference exp(-pi^2/ln lambda)
\newcommand{\NumTangencyRefNear}{\ensuremath{4.6\times10^{-31}}}
% source: e07 tangency k+1=0.025: heuristic splitting reference
\newcommand{\NumTangencyRefMid}{\ensuremath{1.9\times10^{-8}}}
% source: e07 tangency cutoff 1e-12, k+1 <= 0.025: smallest per-config min |sin|
\newcommand{\NumTangencyMinSinFineLo}{\ensuremath{2.0\times10^{-6}}}
% source: e07 tangency cutoff 1e-12, k+1 <= 0.025: largest per-config min |sin|
\newcommand{\NumTangencyMinSinFineHi}{\ensuremath{4.9\times10^{-5}}}
% source: e07 tangency cutoff 1e-12, k+1 <= 0.025: median over configs of the median |sin|
\newcommand{\NumTangencyMedianSinFine}{\ensuremath{4.8\times10^{-4}}}
% source: e07 tangency: WARNING records of any kind, all configs
\newcommand{\NumTangencyWarningsTotal}{0}
% source: e08: solves with a converged orbit (minimal period measured)
\newcommand{\NumSolverConverged}{8\,450}
% source: e08: lower-period solves at requested period 2 with k < 3 (no period-2 orbit)
\newcommand{\NumSolverLowerNoOrbit}{1\,198}
% source: e08: those / converged
\newcommand{\NumSolverLowerNoOrbitRate}{14.2\%}
% source: e08: lower-period solves where the requested period has orbits
\newcommand{\NumSolverLowerSlide}{1\,744}
% source: e08: those / converged
\newcommand{\NumSolverLowerSlideRate}{20.6\%}
% source: e08 k = 10: lower-period / converged
\newcommand{\NumSolverLowerSlideRateKTen}{20.1\%}
% source: e08: lower-period solves that construct_fixed_point accepted as a saddle of the requested period
\newcommand{\NumSolverLowerAccepted}{1\,424}
% source: e08: median solve time (ms)
\newcommand{\NumSolverMedianMs}{0.67}
% source: e08: p99 solve time (ms)
\newcommand{\NumSolverPNinetyNineMs}{7.2}
% source: e09: configs where the brute force beat the library
\newcommand{\NumIsectBruteFasterConfigs}{23}
% source: e09: segments of the largest config
\newcommand{\NumIsectLargestSegments}{\ensuremath{4.1\times10^{5}}}
% source: e09 largest config: library time (s)
\newcommand{\NumIsectLargestLibraryS}{6.6}
% source: e09 largest config: brute-force time (s)
\newcommand{\NumIsectLargestBruteS}{9.0}
% source: e09: configs with at least one u x u crossing
\newcommand{\NumIsectUUConfigs}{2}
% source: e09: area_cutoff of every config with u x u crossings
\newcommand{\NumIsectUUCutoff}{\ensuremath{10^{-5}}}
% source: e10: median summed stage time, CPU (s)
\newcommand{\NumDetTotalCpuS}{0.112}
% source: e10: median summed stage time, GPU (s)
\newcommand{\NumDetTotalGpuS}{0.132}
% source: e10: CPU / GPU median summed stage time
\newcommand{\NumDetGpuSpeedup}{\ensuremath{0.85\times}}
% source: e01_map_throughput.jsonl: records
\newcommand{\NumConfigsEOne}{8}
% source: e02_gpu_end_to_end.jsonl: records
\newcommand{\NumConfigsETwo}{33}
% source: e03_scaling.jsonl: records
\newcommand{\NumConfigsEThree}{22}
% source: e04_area_cutoff.jsonl: records
\newcommand{\NumConfigsEFour}{28}
% source: e05_invariants.jsonl: records
\newcommand{\NumConfigsEFive}{8}
% source: e06_parameter_sweep.jsonl: records
\newcommand{\NumConfigsESix}{98}
% source: e07_breaking.jsonl: records
\newcommand{\NumConfigsESeven}{92}
% source: e08_fixed_point_solver.jsonl: records
\newcommand{\NumConfigsEEight}{24}
% source: e09_intersections.jsonl: records
\newcommand{\NumConfigsENine}{26}
% source: e10_determinism.jsonl: records
\newcommand{\NumConfigsETen}{10}
% source: e11_blast_depth.jsonl: records
\newcommand{\NumConfigsEEleven}{53}
% source: all full-run JSONL files (e01-e11): records
\newcommand{\NumConfigsTotal}{402}
% source: e01-e10 full-run JSONL files: records of the scripted campaign (e11 ran on its own)
\newcommand{\NumConfigsCampaign}{349}
% source: progress.log: wall-clock of the last --quick campaign (min)
\newcommand{\NumCampaignQuickMin}{1.2}
% source: progress.log: wall-clock of the full campaign e01-e10 (min; e11 ran on its own, see NumDeepWallMin)
\newcommand{\NumCampaignFullMin}{59.8}
% source: all full-run JSONL files: outcome ok
\newcommand{\NumOutcomeTotalOk}{367}
% source: all full-run JSONL files: outcome exception
\newcommand{\NumOutcomeTotalException}{24}
% source: all full-run JSONL files: outcome oom
\newcommand{\NumOutcomeTotalOom}{11}
% source: all full-run JSONL files: outcome timeout
\newcommand{\NumOutcomeTotalTimeout}{0}
% source: all full-run JSONL files: outcome invariant_assert
\newcommand{\NumOutcomeTotalInvariantAssert}{0}
